\section{Related Work}
\label{sec:related}

\subsection{Reasoning about cloud access control}

Zelkova encodes AWS policies into SMT and answers questions about them
exactly, such as whether one policy is more permissive than
another~\cite{backes2018zelkova}. What it answers is a question about a policy
or a pair of policies. The property here belongs to a deployment: it needs the
principals, the agent runtimes, and the arguments that bind one to another.
Asking whether role A's policy is more permissive than role B's does not
reveal that A can replace what B runs.

\subsection{Multi-step attacks and escalation paths}

Shevrin and Margalit use model checking to find multi-step IAM attacks, where
each step is an API call that changes what the attacker can do
next~\cite{shevrin2023multistep}. TAC combines static and dynamic information
to detect IAM privilege escalation~\cite{hu2023tac}. Both treat escalation as
a path rather than a single bad grant, which is the stance here.

The difference is what sits at the end of an edge. There the attacker moves
between IAM principals: it assumes a role, creates a key, attaches a policy.
Those edges exist here too. The edge this paper adds does not move between
principals at all: the attacker rewrites the \emph{workload} an existing
principal runs, and nothing about that principal changes. An analysis built on
principal to principal transitions has no place to put that edge.

A second difference is the input: model checking needs a deployed account,
while \tool runs against the source, which is what makes the measurement in
Section~\ref{sec:results} possible and puts a finding in review.

\subsection{Reporting on this platform}
\label{sec:industry}

Cross-agent escalation on \agentcore is not an unreported problem, and the
term is not ours. Unit~42 disclosed that the execution role the AgentCore
starter toolkit creates carries wildcard scopes, so a compromised agent can
read other agents' memory, pull their container images, and invoke
them~\cite{hadad2026godmode}. Mitigant built four paths on that same role,
including a cross-agent runtime hijack~\cite{torkura2026agentsore}. Rehberger describes cross-agent escalation
in a different setting, where one agent approves actions for
another~\cite{rehberger2025freeeachother}.

BeyondTrust is the closest: \iam{bedrock-agentcore:\allowbreak Invoke\allowbreak AgentRuntimeCommand}, on
its own, runs a
command as root inside an existing runtime's sandbox, and the platform's
metadata service returns that runtime's execution role
credentials~\cite{garcia2026agentcorepaths}. One permission, no
\iam{iam:PassRole}, no change to the deployment. Its \iam{iam:PassRole} chains
are create variants, for when no runtime carries the role the attacker wants.

\paragraph{Two routes, and the one this paper takes.}
The two answer different questions. Theirs reads a credential and leaves the
agent as it was; a rotation ends it, and the command lands in the log group
the platform creates for the runtime. Replacing the artifact leaves the agent
answering its own callers with attacker code until someone compares the
deployed artifact against what was meant to be there. Section~\ref{sec:case-study}
measures that difference: with every policy detached from the attacker
principal, the operator's own call to the runtime still returns the planted
code's answer, carrying the victim role's data. The cost of that durability is
a second permission, always \iam{iam:PassRole}.
\api{UpdateAgentRuntime} appears in none of the reporting above.

\paragraph{The gate is a name match.}
The wildcard execution role of the earlier reports is a misconfiguration, and
the fix is to scope it. Section~\ref{sec:case-study} isolates a different
control: with the attacker holding only \iam{BedrockAgentCoreFullAccess}, the
single variable that decides the outcome is whether the victim role's
\emph{name} matches \iam{*BedrockAgentCore*}, which the console's own default
names do. BeyondTrust names that managed policy as over broad and proposes a
service control policy keyed on \iam{aws:PrincipalArn}; the role name
condition inside the policy is not examined.

\paragraph{Prevalence.}
Each of those reports is a single scenario in a controlled account. None
measures how often the preconditions occur, and none reads a deployment before
it exists.

\subsection{Infrastructure as code}

Rahman et al.\ catalogue recurring security smells in IaC scripts and show
they are common in public repositories~\cite{rahman2019sevensins}, and open
scanners implement rule sets of that kind. The unit of analysis differs: a
smell fires on one resource, while the property here is a relation between two
resources in the same file. A deployment where every resource passes every
rule is the deployment this paper is about.

\subsection{Security of LLM agents}

Abdelnabi et al.\ show that content an LLM-integrated application retrieves
can carry instructions the application then
follows~\cite{abdelnabi2023notwhat}, which is why
Section~\ref{sec:threat-model} starts from attacker code running as one
agent. That literature stops at the boundary of the agent's own capabilities;
this paper starts there and asks what the account looks like when the agent is
one of several.
